\documentclass[../main.tex]{subfiles}
\begin{document}
\subsection{Need of Convolution}
Normal layers or feed forward layers do affine transformations to the inputs. A vector is received as input and is multiplied with a matrix to produce an output to which a bias vector is usually added before passing the result through a non-linearity.
But data like images and audio clips have a structure like their arrangement, height and width. Images have channels of different colors and audio clips for e.g. stereo audio has left and right channels. Using these channels we can get more out of such data than the normal feed-forward layer.
\subsection{Discrete Convolution}
During discrete convolution we have an input matrix and a kernel. The kernel slides over the input matrix with some stride and then the area of matrix under the kernel is multiplied by the values of kernel and then added, this result is then stored as the result of that stride.
Once the kernel reaches the end of matrix we have a convoluted layer.
\begin{align*}
  \text{input matrix of size} &= m \times n \\
  \text{kernel matrix size} &= i \times j \\
  \text{stride} &= s \\
  \text{padding} &= p
\end{align*}
  The ouptut of the convolutional operation would be of dimensions.
\begin{align*}
  \text{rows} &= \frac{m-k+2p}{s} + 1 \\
  \text{columns} &= \frac{n-j+2p}{s} + 1
\end{align*}
\subsection{Pooling}
Pooling is just like convolution but instead of linear combination we just we just pass the data under the pooling window or kernel to a pooling function and the output from the function is treated as output of discrete convolution. We can calculate the size of the layer after a pooling operation similarly to how we calculated the size of convolutional layer.
\subsection{Conv By Hand}
\begin{align*}
  \text{The Input:   }
\begin{matrix}
  a_{0,0} & a_{0,1} & a_{0,2} & a_{0,3}\\
  a_{1,0} & a_{1,1} & a_{1,2} & a_{1,3}\\
  a_{2,0} & a_{2,1} & a_{2,2} & a_{2,3}\\
  a_{3,0} & a_{3,1} & a_{3,2} & a_{3,3}\\
\end{matrix}\\ \\
\text{The Kernel:   }
\begin{matrix}
  k_{0,0} & k_{0,1} & k_{0,2}\\
  k_{1,0} & k_{1,1} & k_{1,2}\\
  k_{2,0} & k_{2,1} & k_{2,2}\\
  k_{3,0} & k_{3,1} & k_{3,2}\\
\end{matrix}\\ \\
\text{Output of Normal Convolution:    }
\begin{matrix}
  a & b \\
  c & d \\
\end{matrix}\\ \\
a = (a_{0,0}*k_{0,0})+(a_{0,1}*k_{0,1})+(a_{0,2}*k_{0,2})\\+(a_{1,0}*k_{1,0})+(a_{1,1}*k_{1,1})+(a_{1,2}*k_{1,2})\\+(a_{2,0}*k_{2,0})+(a_{2,1}*k_{2,1})+(a_{2,2}*k_{2,2}) \\ \\
b = (a_{0,1}*k_{0,0})+(a_{0,2}*k_{0,1})+(a_{0,3}*k_{0,2})\\+(a_{1,1}*k_{1,0})+(a_{1,2}*k_{1,1})+(a_{1,3}*k_{1,2})\\+(a_{2,1}*k_{2,0})+(a_{2,2}*k_{2,1})+(a_{2,3}*k_{2,2}) \\ \\
c = (a_{1,0}*k_{0,0})+(a_{1,1}*k_{1,2})+(a_{1,3}*k_{0,2})\\+(a_{2,0}*k_{1,0})+(a_{2,1}*k_{1,1})+(a_{2,2}*k_{1,2})\\+(a_{3,0}*k_{2,0})+(a_{3,1}*k_{2,1})+(a_{3,2}*k_{2,2}) \\ \\
d = (a_{1,1}*k_{0,0})+(a_{1,2}*k_{0,1})+(a_{1,3}*k_{0,2})\\+(a_{2,1}*k_{1,0})+(a_{2,2}*k_{1,1})+(a_{2,3}*k_{1,2})\\+(a_{3,1}*k_{2,0})+(a_{3,2}*k_{2,1})+(a_{3,3}*k_{2,2}) \\ \\
\end{align*}
\begin{align*}
\text{Matrix version of this goes like: }\\
a = [a_0,a_1,a_2,a_3,a_4,a_5,a_6,a_7,a_8,a_9,a_{10},a_{11},a_{12},a_{13},a_{14},a_{15}] \\
k = \begin{matrix}
  w_{0,0} & w_{0,1} & w_{0,2} & 0
\end{matrix}
\end{align*}
\end{document}
